\section{Organizing Hypothesis}
\label{sec:hypothesis}

The organizing hypothesis of this research program is:

\begin{quote}
\textbf{Information is the change in the reachable caused by interaction.}
\end{quote}

This is a hypothesis to be operationalized and tested, not an established law
of physics. It motivates a narrower experimental question: can finite local
interactions produce constraints on reachable futures that persist even when
the local state realizing those constraints is replaced?

The primitive experimental target is therefore the latent counterfactual
structure of lawful successor events, rather than an early analogue of mass,
energy, force, or a classical field. The lattice stores state; a marked
relational transition graph represents which successor interactions remain
possible; and coherence is provisionally the organization that preserves and
transmits part of that structure across realized transitions. Familiar physical
observables, if they emerge, are downstream quantities to explain rather than
primitive variables to imitate. Accordingly, a proposed experiment must state
which aspect of the successor structure it causally changes and measures.

For a state $s_t$, define the locally reachable next-state set as

\begin{equation}
\mathcal{R}_x(t+1\mid s_t)
=
\left\{(A',E'):(A',E')\text{ is admissible under }\mathcal{L}\right\}.
\label{eq:reachable}
\end{equation}

When interaction strictly reduces the reachable set, a cardinality measure is

\begin{equation}
\Delta I_x
=
\log\frac{|\mathcal{R}_x^{\mathrm{prior}}(t+1)|}
{|\mathcal{R}_x^{\mathrm{posterior}}(t+1)|}.
\label{eq:information}
\end{equation}

Equation~\ref{eq:information} does not distinguish all set reshapings that
preserve cardinality. A set- or distribution-sensitive extension remains open.

\subsection{Relation to prior work}

The logarithm of a finite number of alternatives is Hartley information
\cite{hartley1928transmission}; therefore Equation~\ref{eq:information} is a
difference of Hartley support measures when its alternatives are treated
uniformly, not a new information measure. Restriction-centered accounts have
also identified information with limitations on possible variable values
\cite{zadeh2015information}. Cause--effect repertoires and effective information
likewise quantify how mechanisms or interventions constrain possible states in
discrete dynamical systems \cite{balduzzi2008integrated,hoel2013causal}.

Counterfactual approaches to semantic information test whether scrambling or
erasure changes a system's future or viability
\cite{kolchinsky2018semantic,bartlett2024causal,sowinski2025causal}. Persistence
of organization through component turnover is central to autopoiesis and
constraint-closure accounts \cite{varela1974autopoiesis,beer2004autopoiesis,
montevil2015closure}. Related operational tools include local causal states for
coherent structures, reachability in temporal networks, and causal ancestry in
cellular automata \cite{rupe2018local,holme2012temporal,hintze2026rethinking}.

At the physical-foundations level, constructor theory formulates information in
terms of possible and impossible transformations
\cite{deutsch2015constructor}. A recent preprint also describes information and
computation through persistent relational constraints on
admissible transitions \cite{witte2026computation}. The present work therefore
does not claim novelty for information-as-restriction, organization across
turnover, or causal intervention on future possibilities. Its narrower novelty
target is the exact operational combination of disjoint bridge replacement,
paired stochastic intervention, and ordered realization cuts in this finite
local system. Whether that diagnostic is new and general remains to be tested
against the cited methods and matched controls.

The experiments also distinguish a state-level possibility from its lifetime
under ordered resolution. Let $\mathcal{R}^{(k)}$ denote the routes still
admissible after event position $k$ within a declared update. A route may be
lawful in $\mathcal{R}^{(k)}$ yet disappear from $\mathcal{R}^{(k+1)}$ when a
competing lawful event exhausts its source or preempts its target. This notation
does not assert non-Markovian physics: the complete simulator state, RNG state,
and scheduler determine the update. It makes the causal question operational:
must an ancestry-preserving successor be realized before its route closes?
